\documentclass[12pt,a4paper]{article}
\usepackage[utf8]{inputenc}
\usepackage[T1]{fontenc}
\usepackage[margin=2.5cm]{geometry}
\usepackage{graphicx}
\usepackage{listings}
\usepackage{xcolor}
\usepackage{caption}

\newcommand{\studentgroup}{FAF-242}
\definecolor{termbg}{rgb}{0.96,0.96,0.96}
\lstset{backgroundcolor=\color{termbg},basicstyle=\ttfamily\footnotesize,breaklines=true,frame=single}

\begin{document}
\begin{titlepage}
    \centering
    \IfFileExists{utm_logo.png}{\includegraphics[width=0.25\textwidth]{utm_logo.png}\\[0.6cm]}{\vspace*{2cm}}
    {\large \textbf{UNIVERSITATEA TEHNICĂ A MOLDOVEI}}\\[0.2cm]
    {\large Facultatea Calculatoare, Informatică și Microelectronică}\\[0.2cm]
    {\large Departamentul Ingineria Software și Automatică}\\[2.5cm]
    {\Large \textbf{LUCRAREA DE LABORATOR NR. 6}}\\[0.4cm]
    {\LARGE \textbf{Interiorul sistemului de fișiere xv6}}\\[0.4cm]
    {\large Disciplina: Sisteme de Operare}\\[2.5cm]
    \begin{minipage}{0.45\textwidth}\flushleft\textbf{Done by:}\\st. gr. \studentgroup\\Daniel-Ioan \textsc{Strunga}\end{minipage}\hfill
    \begin{minipage}{0.50\textwidth}\flushright\textbf{Checked by:}\\asist. univ. Patricia \textsc{Reițman}\\asist. univ. Marius \textsc{Dumitrașcu}\end{minipage}
    \vfill{\large Chișinău --- 2026}
\end{titlepage}
\newpage

\section{Scopul lucrării}
Studierea structurilor interne de date folosite de un sistem de operare pentru gestionarea memoriei persistente, prin analizarea codului sursă xv6 și crearea comenzilor sistem \texttt{mystat} și \texttt{find}.

\section{Desfășurarea lucrării}
\begin{figure}[h!]
\begin{lstlisting}[language=bash]
$ mystat README 
README: file, inode 2, links 1, size 2441 bytes
$ find / .
/README
/cat
/echo
\end{lstlisting}
\caption{Figure 6.1 -- Rularea comenzilor în mediul virtual xv6}
\end{figure}

\subsection{Răspunsuri (Observe)}
\begin{enumerate}
    \item Un director conține un șir de structuri \texttt{dirent}, fiecare având 16 octeți (2 pentru indexul de inod, 14 octeți max. pentru nume \texttt{DIRSIZ}).
    \item Legăturile hard funcționează iterând atributul \texttt{ip->nlink} din nucleu. Un \texttt{inode} se distruge doar când contorul \texttt{nlink} devine 0 \textit{și} nu mai există descriptori de fișier (\texttt{ref == 0}) care să îl țină deschis.
    \item Limita de mărime a unui fișier în xv6 este 268KB, calculată prin cele 12 blocuri de adresare directă + cele 256 de adrese din blocul de adresare indirectă, înmulțite cu dimensiunea unui bloc.
\end{enumerate}

\section{Concluzii}
Interacțiunea directă cu structurile \texttt{stat} și \texttt{dirent} în limbajul C demonstrează rigiditatea dar și eficiența mecanismelor interne din kernel. Am observat cum abstracțiile precum parcurgerea iterativă a directoarelor (comanda \texttt{find}) se bazează pe citirea și interpretarea secvențială a blocurilor de memorie organizate de xv6.
\end{document}
